\documentclass[10pt,a4paper]{amsart}
\usepackage[margin=19mm]{geometry}
\usepackage{amsmath,amssymb,mathtools}
\usepackage[scr=rsfs,cal=euler]{mathalfa}
\usepackage{xfrac}
\usepackage[unicode,hidelinks,bookmarks=false]{hyperref}
\newcommand{\ie}{i.e.\ }
\newcommand{\cf}{cf.\ }
\newcommand{\resp}{respectively\ }
\newcommand{\Subsection}{\subsection}
\providecommand{\nobreakdash}{\nobreak\hbox{-}}
\setlength{\emergencystretch}{2em}
\multlinegap0pt
\usepackage{xcolor}
\usepackage{amssymb}
\usepackage{tikz}
\newtheorem{lem}{Lemma}
\title{Extremality of principal quiver Grassmannians}
\author{Giovanni Cerulli Irelli, Evgeny Feigin, Markus Reineke}
\date{Source: arXiv:2607.20336v1 (CC BY 4.0)}
\begin{document}
\maketitle

\noindent\emph{Source and licence.} Extremality of principal quiver Grassmannians. Version arXiv:2607.20336v1 (CC BY 4.0), distributed by arXiv under the Creative Commons Attribution 4.0 International licence, http://creativecommons.org/licenses/by/4.0/, as recorded in the arXiv licence metadata for this exact version and shown on the abstract page https://arxiv.org/abs/2607.20336v1. Full licence notice: \texttt{licenses/arxiv-2607.20336-CC-BY-4.0.txt}.\par\smallskip

\noindent\textit{Editorial scope.} Counted unit: the article's lem lem (source0.tex lines 575--578). The article's complete original proof of it is delivered with it (source0.tex lines 580--633, one \texttt{proof} environment of 54 lines). No mathematics is written, repaired or re-derived: the statement and the proof are the article's own lines, in the article's own notation, and no retained line is substituted or re-flowed.  Retained uncounted as the source-local context the proof needs: lines 407--411; lines 435--437.  Nothing was omitted from any retained span, so the package carries no omission record.  No retained line contains a citation, so the wrapper carries no bibliography and no bibliography entry.  No retained line contains a figure environment, an image-inclusion command or drawing code, so no image is delivered and the package declares no asset dependency.  Editorial formatting only, in addition: an AMS \texttt{amsart} preamble that restates the article's own declarations and macros used by the retained lines, namely the environments \texttt{lem} on separate counters, together with the standard packages those lines need.  The retained environments print in this document as Lemma 1 in order of appearance, whereas the article prints the same items with the numbering of its own full document.  The complete native source of the article as distributed in the arXiv e-print for this exact version is bundled unchanged as \texttt{sources/205805/source0.tex}.
\begin{lem}\label{lemij} Suppose that $U\in\mathcal{C}_{\rm simp}$, that $P_i$ is a summand of 
$P_0(U)$, and that $P_j$ is a summand of $P_1(U)$. Then there exists an arrow $i\rightarrow j$, 
and the multiplicity of $P_j$ in $P_1(U)$ equals $1$. 
In particular, if $[U,S_j]^1\not=0$, then we have $[U,S_j]^1=1$.
\end{lem}

\begin{lem}\label{lemp0u} If $U\in\mathcal{C}_{\rm simp}$, then $P_0(U)$ is the direct 
sum of all $P_i$ such that $i$ is a source in the support of ${\rm\bf dim}(U)$.
\end{lem}

\begin{lem} Assume that $P$ contains $P_{\rm min}$ as a direct summand. Then, for every $U\in\mathcal{C}_{\rm simp}$, there exists a direct summand $P'(U)$ of $P$ containing $P_0(U)$ as a subrepresentation, 
such that $${\rm Ext}^1(P'(U)/P_0(U),P)=0,$$
and $P'(U)/P_1(U)$ is indecomposable.
\end{lem}

\begin{proof} 
%\textcolor{blue}{As above, I suggest replacing $I,S$ as subsets of $Q_0$ by some other letters}
By Lemma \ref{lemp0u}, $P_0(U)$ is the direct sum of all $P_i$ for $i$ in $Q_0^I$, the set of sources in the support of ${\rm\bf dim}(U)$. We consider the set of vertices $Q_0$ with its partial order defined by the arrows (that is, $i\leq j$ if there exists a path from $i$ to $j$). Let $Q_0^i$ be the set of vertices admitting a path to $i$ for $i\in I$. We claim that there are no paths between any two different $Q_0^i$. Namely, again by Lemma \ref{lemij}, we know that for any two vertices $i,i'\in Q_0^I$ we find a vertex $j$ and arrows $i\rightarrow j\leftarrow i'$, a contradiction to $Q$ being a tree. Using this information, we can define a specific $P'$ as above. Namely, for every $i\in Q_0^I$, let $Q_0^{K,i}$ be the set of $\leq$-maximal vertices $k(i)\in Q_0^i$ such that $P_{k(i)}$ is a direct summand of $P$. It is nonempty by assumption on $P$. Then we define $P'(U)$ as the direct sum of all $P_k$ for $k$ in $Q_0^K=\bigcup_{i\in I}Q_0^{K,i}$. Since $P_t$ embeds into $P_u$ if $u\leq t$, we see that $P_0(U)$ embeds into $P'(U)$. By definition, $P'(U)$ is a direct summand of $P$. Now we define $R=P'(U)/P_0(U)$, and we claim that $[R,P]^1=0$. Namely, assume that $[R,P_l]^1\not=0$ and that $P_l$ is a direct summand of $P$. From the exact sequence
$$0\rightarrow P_0(U)\rightarrow P'(U)\rightarrow R\rightarrow 0$$
we get an exact sequence
$$
0={\rm Hom}(R,P_l)\rightarrow{\rm Hom}(P'(U),P_l)\rightarrow{\rm Hom}(P_0(U),P_l)\rightarrow{\rm Ext}^1(R,P_l)\not=0.
$$
Thus $[P_0(U),P_l]>[P'(U),P_l]$ which, by the above explicit description of $P_0(U)$ and $P'$, 
means that the number of paths from $l$ to a vertex in $Q_0^I$ is strictly bigger than the number 
of paths from $l$ to $Q_0^K$. Since there are no paths between the sets $Q_0^i$ for $i\in Q_0^I$, 
this implies that, for some $i\in Q_0^I$, there exists a path from $l$ to $i$, but no path from $l$ 
to any vertex in $Q_0^{K,i}$. But since $P_l$ is a direct summand of $P$, the definition of $Q_0^{K,i}$ implies that $l\leq k$ for some vertex in $Q_0^{K,i}$, a contradiction.

Now we consider the following commutative diagram with exact rows and columns:

$$\begin{array}{ccccccc}
&0&&0&&&\\ &\downarrow&&\downarrow&&&\\ &P_1(U)&=&P_1(U)&&&\\
&\downarrow&&\downarrow&&&\\
0\rightarrow&P_0(U)&\rightarrow&P'(U)&\rightarrow&R&\rightarrow 0\\
&\downarrow&&\downarrow&&||&\\
0\rightarrow&U&\rightarrow&V&\rightarrow&R&\rightarrow 0\\
&\downarrow&&\downarrow&&&\\ &0&&0&&&\end{array}$$
Our aim is to prove indecomposability of $V$. Since $[R,P]^1=0$ and $P'(U)$ is a direct summand of $P$, we have $[R,P'(U)]^1=0$, and thus $[R,V]^1=0$ by the right vertical exact sequence. We claim that we also have $[R,V]=0$. Recall from the previous lemma and Lemma \ref{lemij} that 
$$P'(U)=\bigoplus_{k\in Q_0^K}P_k,\; P_0(U)=\bigoplus_{i\in Q_0^I}P_i,\; P_1(U)=\bigoplus_{j\in Q_0^J}P_j,$$
where there are only paths from $Q_0^K$ to $Q_0^I$ and from $Q_0^I$ to $Q_0^J$. This implies that 
$$[P'(U),P_1(U)]=0=[P_0(U),P_1(U)].$$
We then find
\begin{eqnarray*}
[R,V]&=&[P'(U),V]-[P_0(U),V]=\\
&=&[P'(U),P'(U)]-[P'(U),P_1(U)]-[P_0(U),P_1(U)]+[P_0(U),P_1(U)]=\\
&=&[P'(U),P'(U)]-[P_0(U),P_1(U)]=\\
&=&[R,P'(U)]=0.\end{eqnarray*}
The last equality follows since $R$ does not admit a non-zero projective summand: such a summand would induce a direct summand of $P'(U)$ which does not intersect $P_0(U)\subset P'(U)$, in contradiction to the construction of $P'(U)$.
Now we claim that we have $[U,R]=0$. Since $U\in\mathcal{C}$, we have a commuting square of embeddings
$$\begin{array}{ccc}
{\rm Hom}(P_0(U),P_0(U))&\rightarrow&{\rm Hom}(P_0(U),P'(U))\\
\downarrow&&\downarrow\\
{\rm Hom}(P_1(U),P_0(U))&\rightarrow&{\rm Hom}(P_1(U),P'(U))\end{array}$$
and the claim is equivalent to this diagram being Cartesian. Using the explicit form of the three projectives, we can make all four spaces of maps explicit as direct sums of spaces of paths
$$\begin{array}{ccc}
\bigoplus_{i\in Q_0^I} \langle i\leadsto i\rangle&\rightarrow&\bigoplus_{i\in Q_0^I} \langle k\leadsto i\, |\, k\in Q_0^{K,i}\rangle\\
\downarrow&&\downarrow\\
\bigoplus_{i\in Q_0^I}\langle i\leadsto j\, |\, j\in Q_0^J\rangle&\rightarrow&\bigoplus_{i\in Q_0^I}\langle k\leadsto i\mapsto j\, |\, k\in Q_0^{K,i}, j\in Q_0^J\rangle\end{array}$$
Each direct summand corresponding to a vertex $i\in Q_0^I$ constitutes a square of the form
$$\begin{array}{ccc}\mathbb{C}&\rightarrow&Y\\ \downarrow&&\downarrow\\ X&\rightarrow&X\otimes Y\end{array}$$
where the maps are induced by fixed non-zero choices of elements $x\in X$, $y\in Y$. Such a diagram is evidently cartesian, proving the claim $[U,R]=0$. Now from the induced exact sequences
$$0={\rm Hom}(R,V)\rightarrow{\rm Hom}(V,V)\rightarrow{\rm Hom}(U,V)\rightarrow{\rm Ext}^1(R,V)=0$$
and
$$0\rightarrow{\rm Hom}(U,U)\rightarrow{\rm Hom}(U,V)\rightarrow{\rm Hom}(U,R)=0$$
we conclude
$$[V,V]=[U,V]=[U,U]=1$$
since $U$ is indecomposable, thus, $V$ is indecomposable.
\end{proof}

\end{document}
